\section{ZK-GBDT}\label{sec:zkgbdt-protocol}

As outlined in \Cref{sec:technical-overview},
    we commit the forest by committing
    \begin{align}
        \mathcal{F} = (\mathbf{lc}, \mathbf{rc}, \mathbf{N}^{\leftarrow}, \mathbf{N}^{\rightarrow}, \mathbf{E}, \boldsymbol{\Theta}, \mathbf{tree}, \mathbf{root}, \mathbf{score}, \mathbf{is\_leaf}).
    \end{align}

\subsection{Proof Primitives}
\paragraph{Zero-check.}
Given $f(X)$ encoding vector $\mathbf{f}$ over $\mathbb{H}_n$,
    the zero-check validates that $\mathbf{f} = \mathbf{0}$ by proving the existence of $q(X)$ such that
    \begin{align}
        f(X) = q(X)\cdot \nu_n(X).
    \end{align}

We note that as noted in \Cref{thm:domain-lifting-batching}, all zero-checks can be batched together.
    
\paragraph{Sum-check.}
Given $f(X)$ encoding vector $\mathbf{f}$ over $\mathbb{H}_n$,
    the Aurora sum-check~\cite{aurora19} validates that $\sum_{i=0}^{n-1} \mathbf{f}[i] = s$ by proving the existence of $q(X)$
    and $r(X)$ such that
    \begin{align}
        \deg r(X) < n-1, \\
        f(X) - s/n = r(X) \cdot X + q(X)\cdot \nu_n(X).
    \end{align}
    
We note that as noted in \Cref{thm:domain-lifting-batching}, all sum-checks can be batched together.

\paragraph{Table lookup.}
Given $(\mathbf{F}_i)_{i\in [k]}$ and $(\mathbf{T}_i)_{i\in [k]}$
    encoding matrices $\mathbf{F}_i\in \mathbb{F}^{n\times m_i}$
    and $\mathbf{T}_i\in \mathbb{F}^{N\times m_i}$,
    the lookup relation validates that there exist
    $\mathbf{j}\in [N]^n$, such that
    \[
        \mathbf{F}_i[r, c] = \mathbf{T}_i[\mathbf{j}[r], c], \forall r\in [n], c\in [m_i], i\in [k].
    \]
    This is reduced to a logup instance~\cite{logup22},
    which is further batched together by 
    interleaving batching, c.f. \Cref{lem:interleaving}.
    See \textcolor{red}{Appendix}.

\paragraph{Quantization.}
Given $f(X)$ and $g(X)$ encoding vectors $\mathbf{f}$ and $\mathbf{g}$ over $\mathbb{H}_n$,
    and given a public quantization factor $q$,
    the quantization relation validates that
    \begin{align}
        \mathbf{f}[i] = \left\lfloor \frac{\mathbf{g}[i]}{2^q} \right\rfloor, \forall i\in [n].
    \end{align}
    This is reduced to proving the existence of $\mathbf{q}$ and $\mathbf{r}$ such that
    \begin{align}
        \mathbf{g}[i] = 2^q \cdot \mathbf{f}[i] + \mathbf{r}[i], \forall i\in [n], \\
        0 \leq \mathbf{r}[i] < 2^q, \forall i\in [n].
    \end{align}

As discussed earlier, we first batch all quantization checks
    by \Cref{lem:interleaving}.
    Then proving the range of $\mathbf{r}$ with lookup argument.

\paragraph{Prefix sum.}
Given $f(X)$ encoding vector $\mathbf{f}$ over $\mathbb{H}_{abc}$,
    we require three types of prefix sums:
    prove $\mathbf{g}$ and $\mathbf{h}$ such that
    \begin{enumerate}
        \item $\mathbf{g}[i, j, k] = \sum_{k'=0}^{k} \mathbf{f}[i, j, k']$ and $\mathbf{h}[i, j] = \sum_{k'=0}^{c} \mathbf{f}[i, j, k']$;
        \item $\mathbf{g}[i, j, k] = \sum_{j'=0}^{j} \mathbf{f}[i, j', k]$ and $\mathbf{h}[i, k] = \sum_{j'=0}^{b} \mathbf{f}[i, j', k]$;
        \item $\mathbf{g}[i, j, k] = \sum_{i'=0}^{i} \mathbf{f}[i', j, k]$ and $\mathbf{h}[j, k] = \sum_{i'=0}^{a} \mathbf{f}[i', j, k]$.
    \end{enumerate}
    We note that higher dimensional prefix sums can be reduced to 
    one of the three types,
    e.g. the prefix sum along the second dimension of
    $\mathbf{V}\in \mathbb{F}^{a\times b\times c\times d}$
    is reduced to proving the second type over
    $\mathbf{V}_{\{a\times b\times cd\}}$.
    
For brevity, we defer the concrete constructions to \textcolor{red}{Appendix}.

\paragraph{Column-replication.}
Given $f(X)$ encoding $\mathbf{f}\in \mathbb{F}^{n}$
    and $G(X)$ encoding $\mathbf{G}\in \mathbb{F}^{n\times m}$,
    the column-replication relation validates that 
    \begin{align}
        \forall i\in [n], j\in [m].\ \mathbf{G}[i][j] = \mathbf{f}[i].
    \end{align}
    We defer the concrete construction to \textcolor{red}{Appendix}.

\subsection{Proof of Well-formedness}

Following \Cref{subsec:proof of well-formedness},
    we prove the well-formedness of the committed forest by proving the following properties.

\paragraph{Tree index consistency.}
The prover commits $\mathbf{tree\_lc}$ and $\mathbf{tree\_rc}$ such that
    \begin{align}
        \mathbf{tree\_lc}[i] = \mathbf{tree}[\mathbf{lc}[i]], \\
        \mathbf{tree\_rc}[i] = \mathbf{tree}[\mathbf{rc}[i]],
    \end{align}
    and prove their correctness by proving the lookup relation
    \begin{align}
        \mathbf{lc}\| \mathbf{tree\_lc} \subseteq (0, 1, \dots, N_{\mathrm{tot}}-1)^\top \| \mathbf{tree},\\
        \mathbf{rc}\| \mathbf{tree\_rc} \subseteq (0, 1, \dots, N_{\mathrm{tot}}-1)^\top \| \mathbf{tree}.
    \end{align}
    The prover then uses zero-check over $\mathbb{H}_{N_{\mathrm{tot}}}$ to prove
    \begin{align}
        (\mathbf{1}-\mathbf{is\_leaf})\odot (\mathbf{tree\_lc}-\mathbf{tree}) = \mathbf{0}, \\
        (\mathbf{1}-\mathbf{is\_leaf})\odot (\mathbf{tree\_rc}-\mathbf{tree}) = \mathbf{0}.
    \end{align}

\paragraph{Consistency of leaf nodes.}
To prove that leaves have no child, the prover uses zero-check over $\mathbb{H}_{N_{\mathrm{tot}}}$ to validate
    \begin{align}
        \mathbf{is\_leaf} \odot \mathbf{lc} = \mathbf{0}, \\
        \mathbf{is\_leaf} \odot \mathbf{rc} = \mathbf{0}.
    \end{align}

\paragraph{Hyper-box of roots.}
To prove that the hyper-box of any root is full, the prover uses zero-check over $\mathbb{H}_{N_{\mathrm{tot}}d}$ to validate
    \begin{align}
        \mathbf{Root} \odot \mathbf{N}^{\leftarrow} = \mathbf{0}, \\
        \mathbf{Root} \odot \bigl(\mathbf{N}^{\rightarrow} - (B + 1)\mathbf{1}\bigr) = \mathbf{0}.
    \end{align}

\paragraph{Hyper-box consistency of children.}
To prove that the child hyper-boxes are consistent with the parent and the selected split,
    the prover first commits
    $\mathbf{N}^{\leftarrow}_{\mathrm{lc}}, \mathbf{N}^{\rightarrow}_{\mathrm{lc}},
    \mathbf{N}^{\leftarrow}_{\mathrm{rc}}, \mathbf{N}^{\rightarrow}_{\mathrm{rc}}
    \in \mathbb{F}^{N_{\mathrm{tot}}\times d}$
    and proves that
    \begin{align}
        \mathbf{lc}\| \mathbf{N}^{\leftarrow}_{\mathrm{lc}}\| \mathbf{N}^{\rightarrow}_{\mathrm{lc}}
        &\subseteq
        (0, 1, \dots, N_{\mathrm{tot}}-1)^\top\| \mathbf{N}^{\leftarrow}\| \mathbf{N}^{\rightarrow},\\
        \mathbf{rc}\| \mathbf{N}^{\leftarrow}_{\mathrm{rc}}\| \mathbf{N}^{\rightarrow}_{\mathrm{rc}}
        &\subseteq
        (0, 1, \dots, N_{\mathrm{tot}}-1)^\top\| \mathbf{N}^{\leftarrow}\| \mathbf{N}^{\rightarrow}.
    \end{align}
    Let $\mathbf{Int}\in\mathbb{F}^{N_{\mathrm{tot}}\times d}$ be generated by
    column-replication from $\mathbf{1}-\mathbf{is\_leaf}$.
    The prover then uses zero-check over $\mathbb{H}_{N_{\mathrm{tot}}d}$ to validate
    \begin{align}
        \mathbf{Int}\odot\bigl(\mathbf{N}^{\leftarrow}-\mathbf{N}^{\leftarrow}_{\mathrm{lc}}\bigr) &= \mathbf{0},\\
        \mathbf{Int}\odot\bigl(\mathbf{N}^{\rightarrow}-\mathbf{N}^{\rightarrow}_{\mathrm{rc}}\bigr) &= \mathbf{0},\\
        \mathbf{Int}\odot\mathbf{E}\odot
        \bigl(\mathbf{N}^{\rightarrow}_{\mathrm{lc}}-\mathbf{N}^{\leftarrow}_{\mathrm{rc}}\bigr) &= \mathbf{0},\\
        \mathbf{Int}\odot(\mathbf{1}-\mathbf{E})\odot
        \bigl(\mathbf{N}^{\leftarrow}-\mathbf{N}^{\leftarrow}_{\mathrm{rc}}\bigr) &= \mathbf{0},\\
        \mathbf{Int}\odot(\mathbf{1}-\mathbf{E})\odot
        \bigl(\mathbf{N}^{\rightarrow}-\mathbf{N}^{\rightarrow}_{\mathrm{lc}}\bigr) &= \mathbf{0}.
    \end{align}

    Note that the above equations hold for dummy node $0$,
    since all its attributes are zero.


\paragraph{Feature-selection matrix.}
To prove that $\mathbf{E}$ is boolean and one-hot on internal nodes and zero on leaves,
    the prover first uses zero-check over $\mathbb{H}_{N_{\mathrm{tot}}d}$ to validate
    \begin{align}
        \mathbf{E}\odot(\mathbf{E}-\mathbf{1}) = \mathbf{0}.
    \end{align}
    It then invokes the prefix-sum primitive along the feature dimension to obtain
    $\mathbf{E}_{\mathrm{sum}}\in\mathbb{F}^{N_{\mathrm{tot}}}$, where
    \begin{align}
        \mathbf{E}_{\mathrm{sum}}[i]=\sum_{j=1}^{d}\mathbf{E}[i,j].
    \end{align}
    Finally, the prover uses zero-check over $\mathbb{H}_{N_{\mathrm{tot}}}$ to validate
    \begin{align}
        \mathbf{E}_{\mathrm{sum}}+\mathbf{is\_leaf}-\mathbf{1}=\mathbf{0}.
    \end{align}

\paragraph{Scores of internal nodes.}
To prove that internal nodes carry no leaf score, the prover uses zero-check over
    $\mathbb{H}_{N_{\mathrm{tot}}}$ to validate
    \begin{align}
        \mathbf{score}\odot(\mathbf{1}-\mathbf{is\_leaf})=\mathbf{0}.
    \end{align}

\paragraph{Number of roots.}
To prove that the committed forest contains exactly $TK$ trees,
    the prover first uses zero-check over $\mathbb{H}_{N_{\mathrm{tot}}}$ to validate
    the vector root is boolean,
    \begin{align}
        \mathbf{root}\odot(\mathbf{root}-\mathbf{1})=\mathbf{0}.
    \end{align}
    It then uses sum-check over $\mathbb{H}_{N_{\mathrm{tot}}}$ to prove
    \begin{align}
        \sum_{i=0}^{N_{\mathrm{tot}}-1}\mathbf{root}[i]=TK.
    \end{align}
